\section{Representing semantic rules} \label{sec:semantics_lean}

The basic semantics of ShEx formalizes the \lean{Satisfies} ($\vdash$) and \lean{Matches} ($\vDash$)
 relations of Fig. 2 in~\cite{BGP17}, both parameterized by a typing \lean{t}.
The two relations are mutually recursive: a shape expression's neighbourhood
 description invokes \lean{Matches} on the node's outgoing triples, and a
 triple expression's shape reference invokes \lean{Satisfies} indirectly,
 through membership in \lean{t}.
Because each relation calls the other, they are declared together as a
 single \lean{mutual} block, mirroring the paper's joint Fig. 2a/2b
 presentation.

Although \lean{Satisfies} and \lean{Matches} look like a classic mutually
 inductive derivability relation, they are encoded as ordinary
 \lean{Prop}-valued recursive functions rather than a mutual inductive
 family (an inductive family would declare each rule as a constructor of
 the relation, rather than as a case of a function returning that relation).
 This is necessary because the premise in the rule for negation
 uses \lean{Satisfies} negatively, i.e. under a $\neg$: Lean's positivity
 checker requires every recursive occurrence in an inductive definition to
 appear only positively (never to the left of an implication or under a
 $\neg$), since otherwise the definition could be used to derive a
 contradiction. Ordinary recursive functions are not subject to this
 restriction, so defining \lean{Satisfies} as a function sidesteps it.

\lean{Satisfies} and \lean{Matches} are defined case by case, using Lean's
 equational function syntax: each \lean{| .constructor args => body} line
 supplies the definition for one constructor of the input type. Each such
 equation below corresponds to one (or a disjunction of two) rules from the
 paper.

\begin{leanlisting}
mutual
def Satisfies (G : Graph) (t : Typing) (n : RDFNode) : ShExpr → Prop
  | .valueDescr V => InV n V ∨ BNodeInV n V
  | .neigDescr E => Matches G t (neigh G n) E
  | .and S1 S2 => Satisfies G t n S1 ∧ Satisfies G t n S2
  | .or S1 S2 => Satisfies G t n S1 ∨ Satisfies G t n S2
  | .not S => ¬ Satisfies G t n S

def Matches (G : Graph) (t : Typing) (N : Finset Triple) : TExpr → Prop
  | .tPattern P L =>
      ∃ subj pred obj, N = {(⟨subj, pred, obj⟩ : Triple)} ∧ pred ∈ P ∧ t.has obj L
  | .empty => N = ∅
  | .someOf E1 E2 => Matches G t N E1 ∨ Matches G t N E2
  | .eachOf E1 E2 =>
      ∃ N1 N2, DisjointUnion N N1 N2 ∧
               Matches G t N1 E1 ∧
               Matches G t N2 E2
  | .rep E min max =>
      ∃ Ns : List (Finset Triple),
        DisjointUnionList N Ns ∧
        min ≤ Ns.length ∧
        (∀ m', max = some m' → Ns.length ≤ m') ∧
        ∀ Ni ∈ Ns, Matches G t Ni E
end
\end{leanlisting}

\lean{InV} is a proposition that passes if \lean{n} 
 is an RDF Node in \lean{V}, 
 \lean{BNodeInV} passes if \lean{n} 
 is a blank node in \lean{V}, 
 and \lean{neigh} is a function that returns the set of triples in \lean{G}
 that have \lean{n} as subject. \lean{DisjointUnion} and \lean{DisjointUnionList} are defined as follows:

\begin{leanlisting}
def DisjointUnion (N N1 N2 : Finset Triple) : Prop :=
  Disjoint N1 N2 ∧ N = N1 ∪ N2

def DisjointUnionList (N : Finset Triple) 
                      (Ns : List (Finset Triple)) : Prop :=
  Ns.Pairwise Disjoint ∧ N = Ns.foldr (· ∪ ·) ∅
\end{leanlisting}

  

